\documentclass[12pt]{article}
\usepackage{graphicx} % Required for inserting images
\usepackage{hyperref}
\usepackage{makecell}
\usepackage{makeidx}
\usepackage{eurosym}
\newcommand{\firstPage}{
    \thispagestyle{empty}
    \begin{figure}
    \centering
    \includegraphics[scale=0.7]{Swellfish_logo.png}
    \end{figure}
    \author{
        \date{}
        \href{mailto://swellfish14@gmail.com}{swellfish14@gmail.com} \\
    } 
}
\usepackage{hyperref}
\usepackage{array}
\usepackage{tabularx}
\usepackage{adjustbox}

\newcounter{verscount}
\setcounter{verscount}{0}
\newcommand{\addversione}[5]{
	\ifdefined\setversione
		\setversione{#1}
	\else\fi
	\stepcounter{verscount}
	\expandafter\newcommand%
		\csname ver\theverscount \endcsname{#1&#2&#3&#4&#5}
}

\newcommand{\listversioni}{
	\ifnum\value{verscount}>1
		\csname ver\theverscount \endcsname
		\addtocounter{verscount}{-1}
		\\\hline
		\listversioni
	\else
		\csname ver\theverscount \endcsname\\\hline
	\fi
}

\newcommand{\makeversioni}{
	\begin{center}
		\begin{tabularx}{\textwidth}{|c|c|X|X|X|}
		\hline
		\textbf{Versione} & \textbf{Data} & \textbf{Redattore} & \textbf{Verificatore} & \textbf{Descrizione} \\
		\hline
		\listversioni
		\end{tabularx}
	\end{center}
	\clearpage
}

\begin{document}

\graphicspath{ {../templates/img/} }

\title{Norme di progetto}

\firstPage

\maketitle

\begin{center}
    \begin{tabular}{r | l}
		\multicolumn{2}{c}{\textit{Informazioni}}\\
		\hline
		
			\textit{Redattori} &
			[Davide Porporati, Elena Marchioro, Francesco Naletto]\makecell{}\\

			\textit{Revisori} &
			[Jude Vensil Braceros]\makecell{}\\
			\textit{Responsabili} &
			[Andrea Veronese]\makecell{}\\
		      \textit{Uso} & 
                [Esterno]\makecell{}\\
    \end{tabular}
\end{center}

\begin{center}
    \textbf{Descrizione}\\
    File contenente l'analisi dei requisiti necessaria per la realizzazione del progetto.
\end{center}

\pagebreak

\tableofcontents
\pagebreak

\printindex 

\addversione{0.0.0}{24/04/2023}{Andrea Veronese}{Davide Porporati}{Creata struttura di base del documento}
\addversione{0.0.1}{24/04/2023}{Davide Porporati, Elena Marchioro, Francesco Naletto}{Jude Vensil Braceros}{Aggiunti casi d'uso di base}
\addversione{0.0.2}{26/04/2023}{Davide Porporati, Elena Marchioro, Francesco Naletto}{Jude Vensil Braceros}{Aggiunta nuovi casi d'uso come da capitolato Lumos Minima}
\addversione{0.0.3}{05/05/2023}{}{Jude Vensil Braceros}{Verifica generale del documento}
\makeversioni

\section{Introduzione}
Lo scopo di questo documento è quello di raccogliere i risultati dell'attività di analisi dei requisti, includendo la descrizione dei casi d'uso del software e i requisiti necessari per la sua realizzazione.
Questa analisi nasce dalla necessità di dimostrare di aver capito a fondo i requisiti del problema e le aspettative della soluzione che il nostro gruppo andrà a proporre.

\subsection{Riferimenti}
Questo documento contiene un analisi tecnica dei requisiti necessari allo sviluppo del software per il capitolato C2 di Imola Informatica, Lumos Minima.

\section{Descrizione del prodotto}
L'azienda Imola Informatica propone lo sviluppo di una "$webapp$" che consenta la gestione automatica o manuale degli impianti di illuminazione pubblica "$smart$", ovvero dotati di sensori in grado di rilevare la presenza di persone o veicoli in transito.
Lo scopo del progetto è quello di fornire una piattaforma che grazie all'automazione permetta di ridurre i consumi elettrici e al contempo aumentare la sicurezza stradale dei luoghi in cui sarà installata.
Il funzionamento di tale sistema avviene secondo delle fasi ben precise:
\begin{itemize}
	\item l'illuminazione viene impostata a un livello standard, ovvero un numero compreso tra 0 (spento) e 10(luminosità massima)
	\item Se un sensore rileva un veicolo o una persona, aumenta la luminosità di tutti i lampioni della zona a un livello prefissato
	\item Viene fatto trascorrere un intervallo temporale predefinito, nell'ordine di alcuni minuti
	\item Se al termine dell'intervallo stabilito non ci sono nuovi rilevamenti da parte dei sensori, la luminosità viene abbassata nuovamente al livello pre-rilevamento.
\end{itemize}
Oltre al funzionamento automatico il capitolato prevede la possibilità di un funzionamento manuale, ovvero deve permettere al gestore del sistema di aumentare o ridurre a piacimento la luminosità in una data area, senza che ci sia un effettivo rilevamento.

\subsection{Parti del prodotto}
Il sistema software sarà composto dalle seguenti unità:
\begin{itemize}
	\item webapp, ovvero la dashboard che permette di gestire la aree illuminate
	\item interfacciamento con i lampioni e con i sensori, che avverrà tramite chiamate api-rest
	\item schermata di login/logout, per consentire l'accesso al sistema solamente ai soggetti autorizzati
	\item tracciamento dei guasti
	\item schermata per aggiungere lampioni/sensori o intere aree illuminate.
\end{itemize}
\subsection{Caratteristiche degli utenti}

\subsection{Vincoli e preferenze}
L'azienda proponente non ha imposto vincoli teconologici, ma ha dato dei suggerimenti da considerare:
\begin{itemize}
	\item Utilizzare framework open-source e ben noti, che non comportino il pagamento di canoni mensili o di licenze
	\item L'interfacciamento con i sensori o con i lampioni potrà avvenire tramite api-rest oppure tramite un broker mqtt come "$Mosquito$"
\end{itemize}

\section{Casi d'uso}
\subsection{Attori}
\begin{itemize}
	\item Amministratore di sistema: persona con apposite credenziali e permessi che può gestire per intero il sistema di illuminazione. Può aggiungere, modificare o eliminare aree di illuminazione
	\item Sensore: dispositivo IoT che rileva la presenza in una determinata area
	\item Lampione: fonte luminosa comandata dal sensore o dall'amministratore in modalità manuale
	\item webapp: dashboard che permette di controllare la luminosità del sistema e di ottenere informazioni sul suo stato.
\end{itemize}
\subsection{User Case 1 - Consultazione manuale Lumos Minima}
Tipologia: obbligatorio
\begin{itemize}
	\item Attori primari: Amministratore di sistema
	\item Precondizioni: il gestore dell'illuminazione usa il sistema e l'amministratore vorrebbe avere una panoramica dettagliata sul suo funzionamento.
	\item Postcondizioni: l'amministratore consulta il manuale d'uso di Lumos Minima
	\item Scenario Principale: Al gestore dell'impianto di illuminaizone viene fornito il software con il manuale d'uso e l'amministratore è libero di consultarlo.
\end{itemize}
\subsection{User Case 2 - Accesso al sistema}
Tipologia: obbligatorio
\begin{itemize}
	\item Attori primari: Nuovi amministratori di sistema
	\item Attori secondari: Imola Informatica
	\item Precondizioni: i nuovi amministratori vorrebbero poter accedere al sistema, ma non sono autorizzati
	\item Postcondizioni: Imola Informatica, dopo essere stata contattata dagli amministratori interessati, fornisce le credenziali per l'accesso al sistema
	\item Scenario Principale: gli amministratori non ancora abilitati all'accesso richiedono le credenziali ad Imola Informatica
\end{itemize}

\subsection{User Case 3 - Visualizzazione stato sistema}
Tipologia: obbligatorio
\begin{itemize}
	\item Attori primari: amministratori di sistema
	\item Attori secondari: webapp
	\item Precondizioni: gli amministratori vogliono avere una panoramica sulla luminosità del sistema e sullo stato operativo delle aree gestite
	\item Postcondizioni: la webapp fornisce queste informazioni
	\item Scenario Principale: gli amministratori vogliono verificare se i sistemi operano correttamente
\end{itemize}

\subsection{User Case 4 - Aumento luminosità automatico}
Tipologia: obbligatorio
\begin{itemize}
	\item Attori secondari: sensori
	\item Precondizioni: uno o più sensori rilevano la presenza di pedoni o veicoli all'interno del loro raggio d'azione e il sistema riceve una segnalazione in modalita "$push$".
	\item Postcondizioni: la luminosità dell'area in cui si trova il sensore viene incrementata
	\item Scenario Principale: uno o più sensori rilevano la presenza e aumentano la luminosità dell'area per un periodo di tempo prefissato
\end{itemize}

\subsection{User Case 5 - Aumento luminosità manuale}
Tipologia: obbligatorio
\begin{itemize}
	\item Attori primari: amministratori di sistema
	\item Attori secondari: lampioni
	\item Precondizioni: gli amministratori vogliono aumentare manualmente la luminosità di una zona
	\item Postcondizioni: tutti i lampioni della zona vedono un incremento della luminosità
	\item Scenario Principale: gli amministratori vogliono aumentare la luminosità di una zona in maniera manuale, senza che ci sia un riscontro di presenza effettivo.
\end{itemize}

\subsection{User Case 6 - Inserimento nuova area illuminata}
Tipologia: obbligatorio
\begin{itemize}
	\item Attori primari: amministratori di sistema
	\item Attori secondari: webapp
	\item Precondizioni: gli amministratori vogliono aggiungere una nuova zona illuminata da gestire tramite la webapp
	\item Postcondizioni: viene aggiunta la nuova zona alla lista di quelle già presenti a sistema ed è possibile effettuare le operazioni consentite dal software.
	\item Scenario Principale: gli amministratori vogliono inserire una nuova zona dotata di lampioni controllati da sensori smart 
\end{itemize}

\subsection{User Case 7 - Rimozione area illuminata}
Tipologia: obbligatorio
\begin{itemize}
	\item Attori primari: amministratori di sistema
	\item Attori secondari: webapp
	\item Precondizioni: gli amministratori vogliono rimuovere una  zona illuminata dalla lista di quelle a sistema
	\item Postcondizioni: viene rimossa la zona dalla lista di quelle già presenti a sistema e non sarà più possibile eseguire le operazioni permesse dal software.
	\item Scenario Principale: gli amministratori vogliono rimuovere una nuova zona già presente a sistema.
\end{itemize}
\subsection{User Case 8 - Modifica informazioni area illuminata}
Tipologia: obbligatorio
\begin{itemize}
	\item Attori primari: amministratori di sistema
	\item Attori secondari: webapp
	\item Precondizioni: gli amministratori vogliono modificare le informazioni di una  zona illuminata da gestire tramite la webapp
	\item Postcondizioni: vengono modificati i dettagli della zona in questione ed è possibile effettuare le operazioni consentite dal software.
	\item Scenario Principale: gli amministratori vogliono modificare le informazioni relative a una zona dotata di lampioni controllati da sensori smart.
\end{itemize}

\subsection{User Case 9 - Inserimento nuovo sensore}
Tipologia: obbligatorio
\begin{itemize}
	\item Attori primari: amministratori di sistema
	\item Attori secondari: webapp
	\item Precondizioni: gli amministratori vogliono inserire un nuovo sensore a sistema per rilevare la presenza in una zona prestabilita. Per compiere quest'operazione è necessario essere in possesso delle seguenti informazioni: 
	\begin{itemize}
		\item Tipo di itnterazione con il sensore: Push o Pull
		\item indirizzo IP
		\item polling time
		\item coordinate geografiche del sensore
		\item raggio d'azione del dispositivo
	\end{itemize}
	\item Postcondizioni: se l'utente ha le credenziali corrette, il sistema accetta l'inserimento di un nuovo dispositivo IoT
	\item Scenario Principale: gli amministratori vogliono inserire uno o più sensori per zona illuminata tramite la webapp. 
	Al termine dell'inserimento il sistema darà un feedback positivo se il sensore è stato inserito senza problemi o un feedback negativo in caso contrario.
\end{itemize}

\subsection{User Case 8 - Rilevamento malfunzionamenti}
Tipologia: opzionale
\begin{itemize}
	\item Attori primari: sensori
	\item Attori secondari: amministratore di sistema
	\item Precondizioni: un sensore rileva una discrepanza tra intensità luminosa impostata e quella rilevata in una zona qualsiasi.
	\item Postcondizioni: l'amministratore segnala la zona di illuminazione come guasta e la inserisce in un apposito elenco e la rimuove da questo elenco solamente quando il malfunzionamento è stato risolto
	\item Scenario Principale: i sensori rilevano un guasto in una zona illuminata e l'amministratore interviene inserendola nell'elenco delle zone con malfunzionamenti, salvo poi rimuoverle quando il problema verrà risolto.
\end{itemize}
\end{document}
